% 本文件由 LaTeX 排版 Agent 根据有效 translation.json 自动生成。
% author=Dionysius the Great; work_id=0632; title=狄约尼削的书信及残篇

\chapter{狄约尼削的书信及残篇}

% structure: p0001 source=h1 target=omitted reason=page-title-already-rendered-by-parent

% structure: p0002 source=h2 target=section reason=relative-heading-rank; base=h2

\section{第一封书信。致\Person{多米提乌斯}与\Person{狄迪穆斯}。}

1. 但要我逐一列出我们那些（殉道者）朋友的名字，实属多余，他们人数众多，且你们并不认识。只需知道其中包括男女，青年与老人，少女与老妇，军人与平民——各个阶层和年龄，其中有些人遭受了鞭打与火刑，有些人则死于刀剑，赢得了胜利，获得了冠冕。然而，对于另一些人而言，即使度过极长的一生也不足以证明他们成为主所悦纳的人；正如在我自己的情况中，直到现在，那足够的时间尚未显现。因此，祂为我保留了另一个合宜的时刻，只有祂自己知道，正如祂所说：\Quote{在悦纳的时候，我应允了你；在救恩的日子，我帮助了你。}\BibleRef{依 49:8}\par

2. 既然你们询问我们所遭遇的事，并希望知道我们境况如何，你们已充分了解了我们的命运：当我们——即加伊乌斯、我本人、福斯图斯、伯多禄和保禄——被百夫长、官吏以及随行的士兵和其他侍从作为囚犯带走时，有一些来自马雷奥蒂斯的人来到我们这里，不顾我们的意愿强拉我们，尽管我们不愿跟随，他们仍强行将我们带走；而加伊乌斯、伯多禄和我本人与其他弟兄分离，被单独关在利比亚一处荒凉贫瘠之地，距离帕拉埃托尼翁有三天的路程。\par

\Emph{稍后他继续写道：}——他们藏在城里，秘密拜访弟兄们。我指的是司铎马克西穆斯、狄奥斯科普斯、德默特琉斯和路济乌斯。因为福斯提努斯和阿奎拉，这些在世上更显赫的人，正在埃及流浪。我还特别指出那些在疾病中幸存下来的执事，即福斯图斯、欧瑟比乌斯和凯勒蒙。关于欧瑟比乌斯，我说他是主从起初就坚固的，并使他胜任热忱地执行对狱中宣信者的服务，以及执行那危险的任务，即料理和埋葬那些已圆满的真福殉道者的遗体。因为直到今日，总督仍以残忍的方式处决一些被带到他那里的，正如我已说过的；同时他用酷刑折磨另一些人，用监禁和锁链消耗其他人，并命令无人可以接近他们，严格检查是否有人被发现这样做。然而，天主仍藉着弟兄们的温柔关怀与热忱，给受压迫者带来安慰。\par

% structure: p0006 source=h2 target=section reason=relative-heading-rank; base=h2

\section{第二封书信。致诺瓦图斯。}

狄约尼削致其弟兄诺瓦图斯，问候。若你如你所说，是违背己意被带走的，那么你自愿退隐将表明事实如此。因为忍受任何苦难原是人之本分，以免分裂天主的教会。为阻止教会分裂而殉道，并不比为拒绝崇拜偶像而殉道更不光荣；依我之见，前者甚至比后者更为高尚。因为在前者，一个人只是为自己的灵魂作证；而在后者，他是为整个教会作证。如今，若你能说服或强迫弟兄们重新同心，你的正直将胜过你的错误；后者将不被归咎于你，而前者将受你称赞。但若你在那些不服的弟兄身上不能成功，则请注意拯救你自己的灵魂。愿你在主内研究和平而安好。\par

% structure: p0008 source=h2 target=section reason=relative-heading-rank; base=h2

\section{第三封书信。致安提约基雅的主教法比乌斯。}

1. 在我们这里，迫害并非始于皇帝敕令，而是比之早一整年。有一个先知兼诗人，是这个城市的敌人，不管他是什么，他此前已激起并煽动外教人群众反对我们，用他们本土迷信的火焰再次点燃他们。受他煽动，并找到行恶的完全自由，他们认为对魔鬼的唯一虔诚和服务，就是屠杀我们。\par

2. 首先，他们抓住一个名叫梅特拉斯的老人，命令他说亵渎的话；他拒绝后，他们用棍棒击打他的身体，用尖利的芦苇撕裂他的脸和眼睛，然后把他拖到郊区，在那里用石头砸死他。接着，他们把一个名叫昆塔的女信徒带到一座偶像庙宇，强迫她朝拜偶像；当她转身不看偶像，并表示厌恶时，他们绑住她的双脚，把她拖过整个城市粗糙的石铺街道，同时把她撞向磨石，鞭打她，直到他们把她带到同一个地方，也在那里用石头砸死她。然后，他们一拥而上，冲进敬畏天主者的房屋，任何虔诚之人他们各自认作邻居的，他们就追赶并将他们带走，抢劫掠夺他们，为自己留下更贵重的财物，而将普通的木制物品散落在路上焚烧，使这些地方呈现出被敌人占领的城市景象。然而弟兄们只是退让离去，像保禄所见证的那些人一样，\BibleRef{希 10:30}他们欣然接受被抢走财物的事。而且据我所知，除了偶然落入他们手中的个别人之外，至今没有人否认过主。\par

3. 但他们也抓了那位最可敬的童贞女\Person{阿波罗尼娅}，她当时已年迈，他们打掉了她所有牙齿，割伤了她的下颌；然后在城前点起大火，威胁要将她活活烧死，除非她肯跟着他们一起说出他们不敬的言词。虽然她似乎暂时求饶，但一被放开，她就急切地跳进火中，被烧死了。他们还在一人家里抓住了一个叫\Person{塞拉皮翁}的人；对他施以残酷的折磨，打断了他所有肢体，然后将他从高楼摔到地上。没有一条路、没有一条通道、甚至没有一条小巷，我们可以行走，无论白天黑夜；因为无论何时何地，他们都一直喊着，若有人拒绝重复他们亵渎的言词，就必须立刻被拖去烧死。这些暴行以这种方式严厉地持续了相当长的时间。但当叛乱和内战适时地降临到这些可怜的人身上时，就将他们的野蛮残暴从我们身上转移开，转向他们自己。我们得以稍事喘息，只要他们没空向我们发泄愤怒。\par

4. 但转眼之间，从那个较为仁慈的统治而来的改变就向我们宣布了；现在威胁要临到我们的恐惧是巨大的。确实，谕令已经到了；它的内容几乎完全与我们主先前所启示的相符，向我们展现了最可怕的恐怖，以至于如果可能的话，连选民也会动摇。所有人都大为惊恐，较有名望的人中，有不少人，而且数量众多，迅速因恐惧而顺从了谕令；其他担任公职的人，因职务本身的必要性而被拉去顺从；另一些人被朋友拖去，被点名后走近不洁不圣的祭献；还有的人苍白颤抖地屈服，仿佛他们不是要献祭，而是要成为偶像的祭品和牺牲，以致被围观的大群人所嘲笑，并向所有人表明他们既懦弱不敢面对死亡，也不敢献祭。但也有其他人以更大的敏捷冲向祭台，坚称他们以前从未是基督徒；我们主的预言性宣告对这些人最为真实，那就是这样的人得救将是困难的。其余的人，有些跟随了\Emph{上述}这些群体之一；有些逃跑了，有些被抓了。其中一些人，\Emph{在持守信德方面}，走到了锁链和监禁的地步；他们中的某些人甚至忍受了数天的监禁，但最终在进入法庭之前就背弃了信德；而另一些人，在忍受了一段时间的酷刑之后，在更进一步的苦难面前退却了。\par

5. 但也有其他的人，是主坚定而蒙福的支柱，他们从主自身得到力量，获得与自己内在信德力量相称且相配的能力和活力，并成为祂王国的可敬见证人。其中第一个是\Person{尤利安努斯}，一个患痛风的人，既不能站立也不能行走，他被安排与另外两个抬着他的人一起。这两个人中，一个立即否认了\Emph{基督}；但另一个，名叫\Person{克罗尼翁}，别名\Person{尤努斯}，与他一起的年迈\Person{尤利安努斯}本人，承认了主，被放在骆驼上穿过全城——你知道，那是一座非常大的城市——并在那高处被鞭打，最后在一场巨大的火焰中被烧死，全体民众围观。一个士兵站在他们被带去行刑的地方，反对民众对他们的无礼侮慢，结果因民众对他的呼喊而被追赶；这位最勇敢的天主士兵，名叫\Person{贝萨}，被逮捕；在为了敬虔的那场伟大战斗中极其英勇地表现后，他被斩首。另一个人，按出生是\Place{利比亚}人，在名字和真实的福佑上都是真\Person{马卡尔}，尽管法官多次试图说服他否认，他没有屈服，因此被活活烧死。接着是\Person{埃皮马库斯}和\Person{亚历山大}，他们被囚禁了很长时间，遭受了无数的痛苦和刮板与鞭笞的折磨，最后也在巨大的火焰中被烧成灰烬。\par

6. 与这些一起还有四个女人。其中有\Person{阿莫纳里乌姆}，一位虔诚的童贞女，她被法官以最无情的方式折磨了很长时间，因为她从一开始就清楚声明，她不会说出任何他命令她重复的话；在她坚守了她的宣认之后，她被带去处决。其他的是最可敬的年迈\Person{默库里亚}和\Person{狄奥尼西娅}，她曾是多位孩子的母亲，却不爱她的孩子胜过爱她的主。当总督羞于再对她们施加无益的折磨，看到自己被妇女打败时，她们死于刀下，没有再受更多的酷刑。因为她们的冠军\Person{阿莫纳里乌姆}确实已经为她们所有人承受了酷刑。\par

7. \Person{赫伦}、\Person{阿特}和\Person{伊西多鲁斯}，他们都是埃及人，与他们一起被交付的还有一个大约十五岁的男孩\Person{狄奥斯库鲁斯}。起初，那\Emph{审判官}以为可以轻易诱骗这少年，就用花言巧语试图欺骗他，但未能得逞；随后又试图用酷刑强迫他，以为这样不难使他屈服，但狄奥斯库鲁斯既没有顺从说服，也没有在恐吓面前退让。其余的人，在身体遭受最残忍的撕裂之后，仍保持了坚忍，审判官就把他们也投入了火焰。但狄奥斯库鲁斯却被释放了，因为他在众人面前表现出的卓越风范，以及在审问时回答的极其智慧的话令审判官惊奇；审判官宣称，因他年幼，给予他更多时间悔改。这位最虔诚的狄奥斯库鲁斯目前与我们同在，等待更伟大的争战和更持久的竞赛。还有一个名叫\Person{内梅西翁}的人，他也是埃及人，被诬告为盗贼的同伙；他在百夫长面前澄清了这项指控，证明这是极不自然的诽谤，但有人告发他是基督徒，他不得不作为囚犯来到总督面前。那极不义的法官对他施加了双倍于盗贼的惩罚，使他遭受酷刑和鞭打，然后将他与盗贼一同投入火中。就这样，这位有福的殉道者按基督的样式得到了荣耀。\par

8. 还有一群士兵，包括\Person{阿蒙}、\Person{芝诺}、\Person{托勒密}和\Person{因格努斯}，以及与他们一起的一位老人\Person{特奥菲卢斯}，他们集体站在审判台前。当某个人作为基督徒受审，并已经倾向于否认时，他们围在四周，咬牙切齿，挤眉弄眼，伸手做出各种身体姿态。当所有人的注意力都转向他们时，还未待任何人抓住他们，他们迅速跑到审判席前，宣称自己是基督徒，这给总督及其同僚带来了恐惧；而那些受审的人，在面临将要受的苦时显得极其勇敢，而审判官们自己却在颤抖。于是，这些人精神高昂地离开审判台，因他们的见证而欢欣鼓舞，因为天主亲自使他们荣耀地得胜。\par

9. 此外，许多其他人被异教徒在城邑和乡村中撕裂。我将举其中一个例子。\Person{伊斯基里翁}为一位统治者充当管事，领取固定工资。他的主人命令这人献祭；他拒绝后，主人辱骂他。他坚持不从，主人就侮辱他；他仍不屈服，主人就拿了一根大棍子，刺穿他的肠子和心脏，杀死了他。我为什么要提及那些被迫在荒野和山中漂泊，被饥饿、干渴、寒冷、疾病、盗贼和野兽夺去生命的人呢？这些幸存者就是他们被拣选和得胜的见证。然而，我要补充一件事作为这些事的例证。有一位名叫\Person{凯勒蒙}的老人，他是\Place{尼罗城}的主教。他与同伴逃往\Place{阿拉伯山}，再也没有回来。弟兄们也多次搜索，未能找到他们的任何踪迹；他们从未找到他们本人或他们的尸体。许多人也被野蛮的\Person{撒拉森人}掳到同一座阿拉伯山为奴。其中一些人经过艰难、只付了巨款才被赎回；另一些人至今未被赎回。兄弟啊，我述说这些事并非没有目的，而是为了让你知道我们遭遇了多么众多和可怕的苦难；这些困苦，也只有最经历的人才能最理解。\par

10. 因此，那些曾与我们同在、如今与基督同坐、共享祂的王国、参与祂的审判、并作祂审判助理的圣洁殉道者，在那里接纳了某些跌倒的、因向偶像献祭而负罪的弟兄。他们看到这些人的悔改可能蒙那不愿罪人丧亡、而愿他悔改的主（\BibleRef{则 33:11}）所悦纳，就证实了他们的真诚，收纳他们，重新团聚，与他们一同聚会，在祈祷和庆节中与他们共融。那么，弟兄们，关于这些人，你们给我们什么建议？我们该怎么办？我们是应该站出来，按那些殉道者所作的决断和判断行事，效法他们的恩慈，以同样仁慈对待他们曾施以怜悯的人？还是应当视他们的决定为不义，自己成为他们在这类事上意见的审判官，使宽恕化为眼泪，推翻既有的秩序？\par

11. 但我要在这里更详细地述说一件在我们中间发生的事：我们中有一位年迈的信徒，名叫塞拉彼翁。他一生无可指摘地度过了漫长岁月，但在试炼（迫害）时期跌倒了。这人多次祈祷（求赦罪），却无人理会他，因为他曾向偶像献祭。他患病后，连续三日不能言语，失去意识。第四日稍有好转，他把孙儿叫到跟前说：\Quote{我儿，你还要耽搁我多久？我恳求你，赶快赦免我。请一位司铎到我这里来。}说完这话，他又不能言语了。那孩子跑去请司铎；但当时是夜间，司铎又患病，因而不能前来。然而，我曾下令，对于临终之人，若当时请求赦免，尤其是先前曾恳切寻求过的，应当予以赦免，好使他们能怀着喜乐的希望离开此世。于是司铎给了孩子一小份圣体，吩咐他用水浸泡后滴入老人的口中。孩子拿着那份圣体回来，走近时，还未进门，塞拉彼翁再次清醒过来，说：\Quote{你来了，我的孩子，司铎不能来；但快照你所受的指示做，让我离去。}孩子将圣体碎片浸入水中，立刻滴入（老人）口中；他咽下一点后，立即断了气。他岂不是很显然被保存了吗？他难道不正是活到得以被赦免，并因他所作诸多善行，藉着罪过的洗除而蒙承认吗？\par

% structure: p0020 source=h2 target=section reason=relative-heading-rank; base=h2

\section{第四封书信。致罗马主教科尔内略。}

除了所有这些，他在收到科尔内略反对诺瓦图的信后，也写信给罗马的科尔内略。在那封信中，他也表明自己曾受赫勒努斯（基里基雅塔尔索的主教）以及与他同在的其他人——即卡帕多细雅的主教菲尔米利安和巴勒斯坦的特奥克提斯托斯——的邀请，在安提约基雅会议上与他们相会，因为那里有些人正试图建立诺瓦图的分裂。此外，他写道，他获悉法比乌斯已去世，德默特里亚努斯被任命为安提约基雅教会主教的继任者。他也写到耶路撒冷的主教，用了这样的措辞：\Quote{真福亚历山大被投入监狱后，安然辞世，进入福乐。}\par

% structure: p0022 source=h2 target=section reason=relative-heading-rank; base=h2

\section{第五封书信。论圣洗圣事的第一封，致罗马主教斯德望。}

然而，我的兄弟，你要知道，所有位于东方以及更远地区的教会，先前处于分裂状态，如今都重新合而为一；各地教会的首领都同心合意，为这出乎意料的和平恢复而极其喜乐。我可以提及安提约基雅的德默特里亚努斯；凯撒勒雅的特奥克提斯托斯；埃利亚的马扎巴内斯（接续已故的亚历山大）；提洛的马里努斯；劳狄刻雅的赫利奥多罗斯（接续已故的特利米德雷斯）；塔尔索的赫勒努斯，以及基里基雅的所有教会；菲尔米利安和整个卡帕多细雅。因为我只举出了较著名的几位主教，以免我的书信过长，或使我的论述对您过于沉重。然而，叙利亚和阿拉伯的全境，对于那里的弟兄们，您不时送去供应，如今又寄去了信件；还有美索不达米亚、本都和比提尼亚；简而言之，各处各地都因这如今建立的同心同德和弟兄友爱而喜乐，并为此光荣天主。\par

% structure: p0024 source=h2 target=section reason=relative-heading-rank; base=h2

\section{同文，另一译本。}

但是，我的兄弟，要知道，整个东方以及更远地区的所有教会，先前分裂的，如今终于回归合一；各地教会首长都同心合意，以难以置信的喜乐为这出乎意料的和平回归而欢欣：即安提约基雅的德默特里亚努斯；凯撒勒雅的特奥克提斯托斯；亚历山大去世后埃利亚的马扎贝内斯；提洛的马里努斯；特利米德雷斯去世后劳狄刻雅的赫利奥多罗斯；塔尔索的赫勒努斯，以及基里基雅的所有教会；菲尔米利安努斯，以及整个卡帕多细雅。我只举出了较著名的几位主教，以免我的信过于冗长，我的论述过于疲劳。确实，整个叙利亚和阿拉伯，您不时向其提供帮助，并且现在已寄出信件；还有美索不达米亚、本都和比提尼亚；总而言之，所有地方都在因这和谐与弟兄之爱而欢欣，赞美天主。\par

% structure: p0026 source=h2 target=section reason=relative-heading-rank; base=h2

\section{第六封书信。致西克斯托主教。}

1. 先前（斯德望）确实曾就赫拉努斯和菲尔米利安努斯，以及所有在基里基雅和卡帕多细雅各地，以及所有邻近省份的人写信，告诉他们因同一缘故，他将离开他们的共融，因为他们为异端者重行圣洗。请思考此事之严重性。因为，据我所知，在最重大的主教会议中，已规定从异端来的人应首先接受\Emph{公教}教义的训练，然后藉圣洗洗净旧而邪恶之酵母的污秽。我屡次写信，在这些事上请他作证并呼吁。\par

稍后，狄约尼削继续写道：——\par

2. 此外，对于我亲爱的同司铎狄约尼削和斐理孟，他们先前赞同斯德望，并就相同的事写信给我，我先前曾以数语写信，但现在更详尽地再次写信。\par

在同一封信中，欧瑟比告诉我们，他通知西克斯托有关\par

撒伯流异端者，当时正盛行，他的话如下：——\par

3. 因为那近期在彭塔波利斯的\Place{托勒迈斯}城所兴起的教义，它恳求并充满对全能天主和我们主耶稣基督之父的亵渎；充满不信和背信弃义，针对祂的独生子及一切受造的首生者，那成为人的圣言，并且它除去对圣神的感知——双方的信函都交给了我，弟兄们也前来讨论它，尽我所能借着天主的恩赐更清楚地陈明——我写了一些信，其副本我已寄给你。\par

% structure: p0033 source=h2 target=section reason=relative-heading-rank; base=h2

\section{第七封书信。致\Person{斐理孟}，司铎。}

我确实专心阅读了书籍并仔细研究异端者的传统，甚至到了以他们可恶的见解腐蚀我灵魂的程度；但我从中也获得了这个益处：我能够在内心驳斥他们，并且比以往更衷心地厌恶他们。当一位司铎级别的弟兄试图劝阻我，担心我会陷入同样的邪恶污秽中，因为他说我的心灵会被玷污（确实如实，如我自己所觉察），但我被来自天主的一个神视所坚固。同时有话语对我说，明确命令我：\Quote{阅读凡来到你手中的一切，因为你适合这样做，你纠正并考验每一个人；并且从它们那里，首先向你显明了相信的原因和机会。}我接受这个神视，认为它符合宗徒的话语，正是这样敦促所有被赋予更大德行的人：\Quote{你们要作精明的兑换银钱者。}\par

然后，\Person{欧瑟比}说，他以插话的方式补充了一些关于所有异端的内容：——\par

这项规则和形式我从我们蒙福的教父\Person{赫拉克拉斯}领受：因为你们，那些来自异端的人，即使他们已从教会堕落，更不用说若他们没有堕落，但当他们被看到参加信友集会，并被指控去听邪道教师的话，且被逐出教会，他在多次祈祷之后，不接纳他们，直到他们公开并当众叙述从他们的对手那里听到的一切。然后他终于接纳他们参加信友集会，绝不要求他们重新领受圣洗。因为他们已经从那次圣洗中预先领受了圣神。\par

再一次，经过充分讨论后，他补充说：——\par

此外我还得知，这习俗并非现在才首次传入非洲人中间；而且很久以前，在以前主教的时代，在许多人口众多的教会中，以及当\Place{依科尼雍}和\Place{西纳德}的弟兄们召开会议时，大多数人也赞成它；我不愿推翻他们的决定，而使他们陷入纷争和争论。因为经上记载：\BibleQuote{不可挪移你邻舍的地界，那是你祖先所设立的。}\BibleRef{申 19:14}\par

% structure: p0039 source=h2 target=section reason=relative-heading-rank; base=h2

\section{第八封书信。致\Person{狄约尼削}。}

因此我们正当地击退\Person{诺瓦替安}，他撕裂了教会，并拉走了一些弟兄走向不敬和亵渎；他提出了一个最不虔敬的关于天主的教义，并诬蔑我们最仁慈的主耶稣基督好像祂不仁慈；除此以外，他还认为神圣圣洗无效并拒绝它，且推翻了信德和宣认，它们位于圣洗之前，并彻底从他们那里赶走了圣神，即使有任何希望存在，要么祂会留在他们内，要么祂会回到他们那里。\par

% structure: p0041 source=h2 target=section reason=relative-heading-rank; base=h2

\section{第九封书信。致\Person{西斯笃二世}。}

实在，弟兄，我需要建议，我恳求你的判断，免得我或许在所遭遇的事上犯错。一位来教会聚会的弟兄，一段时间以来被当作信友，并且在我受祝圣之前，甚至在\Person{赫拉克拉斯}本人担任主教之前，他已参与信友集会；当他参与最近受圣洗者的圣洗时，听到了提问和他们的回答，他来到我这里，流着泪，哀叹自己的命运。他俯伏在我脚前，开始承认并宣称他在异端中所领受的圣洗并非这种样式，也与我们这圣洗毫无共同之处，因为它充满亵渎和不敬。他说他的灵魂被极其痛苦的悲伤刺透，他甚至不敢抬眼向天主，因为他被那些邪恶的言语和事物所启蒙。因此他恳求，藉此最纯洁的圣洗，得以获得义子之恩和恩宠。而我，确实不敢这样做；但我认为长期以来的共融已足以达成此目的。因为我不敢重新更新一位曾聆听感恩经并同他人应答\Quote{阿们}的人，他曾站在祭台前伸手领受祝福的食物并领受了它，且长时间参与我们主耶稣基督的体和血。从此我吩咐他鼓起勇气，以坚定的信德和良好的良心亲近\Emph{圣体}，并分享它们。但他不停地哭泣，不敢接近祭台；尽管被劝勉，他几乎无法忍受参加祈祷。\par

% structure: p0043 source=h2 target=section reason=relative-heading-rank; base=h2

\section{第十封书信。驳\Person{杰尔玛诺}主教。}

1. 如今我在天主面前说话，祂知道我没有撒谎：我逃走既非出于自己的选择，也非没有天主的指示。但在较早的时候，当德西乌斯迫害的\Emph{谕令}确定后，\Person{萨比努斯}在那时刻派了一个\Person{弗鲁门塔里乌斯}来搜捕我。我在屋里待了四天，等待这个\Person{弗鲁门塔里乌斯}的到来。但他却到处搜查其他地方：道路、河流、田野，凡他怀疑我会躲藏或行走之处。他像被一种盲目击打似的，从未找到那屋子；因为他从未想到在追捕之下我还会待在家里。然后，过了四天，天主指示我离开，并以出人意料的方式为我开路，我的家人和我，以及相当多的弟兄，一起出了门。后来发生的事清楚表明这是天主眷顾的安排：其中我们或许也对某些人有所帮助。\par

2. \Emph{然后，在中断之后，他叙述了逃跑后遭遇的事，并补充如下陈述：}——约在日落时，我和与我在一起的人被兵士抓住，被带到了\Place{塔波西里斯}。但藉天主的眷顾，那时\Person{蒂莫泰乌斯}没有与我同在，他也根本没有被捉住。他后来到了那地方，发现屋子空了，有官员看守，而我们被掳为奴。\par

3. \Emph{在另一些事之后，他继续说：}——在他身上，这奇妙的眷顾安排是如何进行的？因为事实将被述说。当\Person{蒂莫泰乌斯}极其慌乱地逃跑时，一个乡下人遇见了他。这人问他为什么匆忙，他如实告诉了他。那人（当时正要去参加某些婚礼；因为他们的习惯是整夜聚在这样的场合）听了这事，便继续走向欢庆的地方，进去向那些坐席的人述说了情况；他们就像听到一个口令似的，不约而同地都站起来，全速涌来。他们冲向我们，大声呼喊；看守我们的兵士立刻逃跑，他们来到我们面前，我们正躺在光秃秃的卧榻上。至于我，天主知道，起初我以为他们是来抢劫我们的强盗；我躺在床架上，只穿着亚麻内衣，便把旁边放着的其余衣服递给他们。但他们叫我起来，尽快离开。于是我明白了他们来的目的，便哭喊、恳求、央求他们走开，别管我们；我乞求他们，若愿对我们行善，就赶在押解我的人之前，砍下我的头。我这样大声呼喊时，那些与我一同经历这一切的同伴知道，他们开始用力扶我起来。我仰面倒在地上；但他们抓住我的手和脚，拖着我，把我抬了出去。目睹这一切的人跟着我——就是\Person{加约}、\Person{法乌斯图斯}、\Person{伯多禄}、\Person{保禄}。这些人也扶起我，急忙将我带出小镇，把我放在没有鞍的驴上，就这样把我带走了。\par

4. 我担心自己，在不得不叙述天主眷顾在我们身上奇妙安排的情况下，有被指控极大愚蠢和麻木的危险。然而，正如有人所说，保守君王的秘密是好的，但揭示天主作为则是光荣的，\BibleRef{多 12:7} 我将直面对\Person{杰曼努斯}的猛烈攻击。我到\Person{埃米利亚努斯}那里并非独自一人；同我一起的还有我的同司铎\Person{马克西姆斯}，以及执事\Person{法乌斯图斯}、\Person{欧瑟比}和\Person{海勒蒙}；还有一位从罗马来的弟兄也与我们同去。\Person{埃米利亚努斯}开头并没有对我说：\Quote{不可举行集会。}那对他而言实在是多余的事，也是一个打算回到首要问题之人最不会做的事；因为他关心的不是我们聚集他人，而是我们自己不做基督徒。因此，他命令我停止这样做，无疑认为若我本人背弃，其他人也会效仿。但我既不无理也不多话地回答他：\Quote{我们必须服从天主，而不是服从人。}\BibleRef{宗 5:29} 此外，我公开见证我只敬拜唯一真天主，不敬拜别的，并且我不能改变那立场，也不能停止做基督徒。于是他命令我们到一个靠近旷野的村庄，名叫\Place{塞弗罗}。\par

5. 也请听我们两人共同说过的、已被记录下来的话。当狄约尼削、福斯托、玛克西穆、玛尔塞禄和赫勒蒙被带到法庭上时，总督埃米利安说：\Quote{我实在以自由之言向你们论及我们君王的仁慈，正如你们所体验的；因为他们给了你们自救的能力，只要你们愿意转向合乎本性的事，敬拜那些也维持他们王国的神祇，并忘记那些违背本性的事。对此你们说什么呢？因为我绝不指望你们会对他们的仁慈忘恩负义，他们的目的确实是引导你们走向更善的道路。}狄约尼削这样回答说：\Quote{并非所有人都敬拜所有的神祇，而是不同的人敬拜他们认为是真神的对象。我们现在敬拜独一的天主，祂是万有的创造者，正是那位将统治权交到他们至圣陛下瓦勒良和加里恩手中的天主。我们敬畏并敬拜祂；我们不断为这些君王的统治祈祷，愿它稳固不变。}埃米利安总督对他们说：\Quote{但如果这位是神，谁阻止你们也敬拜祂，连同那些本性为神的神祇呢？因为你们已受命敬拜神祇，就是那些众所皆知的神祇。}狄约尼削回答说：\Quote{我们不敬拜别的神。}埃米利安总督对他们说：\Quote{我看你们对君王的仁慈既忘恩又麻木。因此你们不得留在此城，而要被遣送到利比亚地区，并定居在一个叫克弗罗的地方；我按照君王的命令选择了此地。你们或任何其他人决不允许举行集会，或进入那些称为墓地的地方。若有人被发现没有前往我命令他去的地方，或被发现参加任何集会，他将自招危险，因为应有的惩罚不会落空。所以，前往你们奉命去的地方吧。}他这样强迫生病的我离开，连一天的宽限也不给。这样我还有什么机会考虑举行或不举行集会呢？\par

6. \Emph{然后，在另一些话之后，他说}：——此外，我们并未因主的临在而退避可见的集会。我却更加热切地努力聚集城中的弟兄，如同我在他们中间一样；因为我身体虽不在，如我所说，心灵却在。在克弗罗，有一个相当大的教会与我们同住，部分是由从城中跟随我们的弟兄组成，部分是从埃及加入我们的。在那里，天主也为我们打开了传道之门。起初我们被逼迫、被石头砸，但过了一段时间，一些外邦人离弃了偶像，归向了天主。因为借着我们，圣言首次在他们中间播种，此前他们从未听过。为显示这正是天主引领我们到他们那里的目的，当我们完成了这一服务后，祂又领我们离开。因为埃米利安似乎有意将我们转移到更艰苦、更像利比亚的地区；他下令将所有的人分散到马勒奥提斯地区，并为每队人分配了乡村。但他特别指示将我们安置在大路旁，以便我们可能首先被捉拿；他显然在安排和计划时，为的是随时决定捉拿我们时能轻易地将我们一网打尽。\par

7. 当我接到命令前往克弗罗时，我对那地方的位置一无所知，甚至此前几乎没听说过它的名字；尽管如此，我还是勇敢而平静地去了。但当我被告知要转移到科鲁提翁地区时，在场的人知道我的感受；因为在此我要自我控诉。起初我确实非常烦恼，感到很不舒服；虽然这些地方碰巧更出名并为我们所熟悉，但人们声称该地区缺乏弟兄，甚至缺乏有品格的人，而且受到旅客骚扰和强盗袭击。然而，当弟兄们提醒我它更靠近城市时，我得到了安慰；尽管克弗罗带给我们与从埃及来的各种弟兄广泛交往，以至于我们得以更大规模地举行神圣集会，但在科鲁提翁，由于城市更近，我们可以更频繁地见到那些真正可爱、与我们关系最密切、最亲爱的人：因为他们会来与我们同住休息，而在更远的郊区，会有不同的和特别的聚会。事情果然如此。\par

8. \Emph{然后，在说了其他一些事后，他再次叙述了所遭遇的事：}— \Person{格尔马努斯}夸耀自己的许多信仰宣认。他当然能说出许多发生在他身上的逆境！但他能否数出自己身上像我们那样多的定罪判决，以及没收财产、公敌宣告、财物掠夺、尊严丧失、世俗荣誉的鄙视、对总督和议员赞美的轻蔑、对敌人威胁的忍耐承受、以及呐喊、危险、迫害、流浪生活、困境的压力和各种磨难，正如我在\Person{德西乌斯}和\Person{萨比努斯}时期所遭遇的，以及我在当前\Person{埃米利亚努斯}的严酷统治下所遭受的？但\Person{格尔马努斯}究竟在哪里出现过？又有何提及？然而，我因\Person{格尔马努斯}而陷入这种巨大的愚妄之中，我退避了；因此，我停止向已经深知这些事的弟兄们详细叙述所有发生的事。\par

% structure: p0052 source=h2 target=section reason=relative-heading-rank; base=h2

\section{第十一封信。致\Person{赫尔马蒙}。}

1. 但\Person{加卢斯}不明白\Person{德西乌斯}的邪恶，也没有预见到是什么导致了他的覆灭。他却在眼前的石头上绊倒了；因为当他的统治处于繁荣局面，事情按他的意愿发展时，他压迫了那些为他的和平与福祉向天主恳求的圣洁之人。因此，迫害他们，他也迫害了为自己所作的祈祷。\par

2. 对\Person{若望}也同样给予了一个启示：\Quote{有一张说大话和亵渎话的口赐给了它，} 他说，\Quote{并且也赐给了它权柄，四十二个月。} 人们会发现\Person{瓦莱里安}的情况令人惊叹；最需要思考的是他在这些事发生之前是多么不同，对天主的人是如此温和与友善。因为在他之前的皇帝中，没有一位像他那样对他们表现出如此仁慈和宽厚的态度；是的，即使是那些被说成公开成为基督徒的人，也没有像他在统治初期那样以极端友善和恩惠接待他们；他的全家那时充满虔诚的人，本身就是一座天主的教会。但埃及术士的首领和会长使他放弃了这种态度，怂恿他杀害和迫害那些纯洁圣洁的人，视他们为那些可咒诅和可憎恶的咒语的对头和障碍。因为确实有，并且一直有一些人，仅凭他们的出现、显现、呼吸和说几句话，就能驱散恶鬼的诡计。但他使他心中萌生了举行不洁的入教仪式、可憎的魔术和可厌的祭献，杀害可怜的儿童，用不幸父亲的子女作祭品，分割新生儿的肠子，并肢解和切割天主所造的生物，仿佛通过这些手段他们能获得幸福。\par

3. \Emph{随后他补充了以下内容：}— 那么，\Person{马克里安}为他们带来了何等辉煌的感恩祭献，为那个他曾寄予希望的帝国！他以前被认为是君主的忠实国库官，却无意于任何合理或有益于公益的事，反而使自己受到那预言的诅咒，即 \Quote{祸哉，那些从自己心里预言，而不顾及公益的人！} 因为他没有分辨那管理万物的眷顾，也没有想到那万有之前、藉着万有、并在万有之上的审判。因此，他也成了天主教会的敌人；此外，他疏远并使自己远离天主的仁慈，尽最大可能逃避祂的救恩。而在这点上，他确实展示了他名字的特殊含义的真实性。\par

4. \Emph{随后，在说了其他一些事后，他接着这样说：}— 因为\Person{瓦莱里安}被这个人怂恿做了这些事，因此遭受了侮辱和责备，正如\Emph{主}藉\Person{依撒意亚}所说的话：\BibleQuote{是的，他们选择了自己的道路，心里喜悦自己的可憎之事；我也要选择他们的戏弄，并惩罚他们的罪恶。} 但这个人（\Person{马克里安}）对帝国的热情使他疯狂，他完全不配，同时由于他残疾的身体无法承担帝国统治的标志，就推出他的两个儿子作为他父亲过犯的承载者。因为在他们身上，天主所说的话显然应验了：\BibleQuote{我必追讨父亲的罪恶于子孙，直到憎恨我的三四代。} 因为他把自己未能满足的邪恶情欲堆在儿子们的头上，从而把自己的邪恶转嫁给他们，也把他们自己对天主的仇恨转移了过去。\par

5. 那人背叛了两位皇帝中的一位，又向另一位开战，很快就连同全家彻底灭亡了。\Person{伽利埃努斯}被拥立，并得到所有人的承认。他既是老皇帝，又是新皇帝；因为他比那些人更早，又比他们活得更久。关于这一点，是\Emph{由主}对\Person{依撒意亚}所说的这样一句话：\BibleQuote{看哪，从起初就有的事已经成就了；还有新事将要兴起。}就像一片乌云遮住太阳，暂时荫蔽它，乌云自己出现在它的位置，但当乌云飘过或消散后，先前升起的太阳再次出现并显示自己：同样，这位\Person{马克里亚努斯}自荐，甚至为自己进入已经确立的\Person{伽利埃努斯}的帝国找到了途径；但现在他不再是\Emph{那个}了，因为他确实从未是过，而另一位，即\Emph{\Person{伽利埃努斯}}，仍如从前一样。他的帝国，如同摆脱了衰老，并净化了以前附着的邪恶，现在变得更加强盛繁荣，现在在更远的地方被看见和听闻，并向四面八方扩展。\par

\Emph{接着，他进一步指出他写下这段记述的确切时间，如下：}我又想起要回顾帝国纪年的日子。因为我看那些最不虔敬的人，他们的名字曾经一时显赫，却在短时间内变得无名。但我们更虔诚、更敬畏天主的君主已经过了他的第七年，现在正是第九年，我们要在这年庆祝节日。\par

% structure: p0059 source=h2 target=section reason=relative-heading-rank; base=h2

\section{书信十二。致\Place{亚历山大城}人。}

1. 对其他人来说，目前的状况似乎并非举行节日的合适时机；这对他们来说当然不是节期，事实上，其他任何时间对他们来说也不是。我这样说，不仅指那些明显悲伤的场合，甚至指人们可能认为最欢乐的所有场合。现在确实一切都转为哀悼，所有人都悲伤，城中充满哀哭声，因为每天都有众多死者。正如在\Place{埃及}长子的事上记载的那样，如今也传来极大的哀号。\BibleQuote{因为没有一个家里没有死者。}但愿仅此而已！\par

2. 诚然，在此之前我们也遭受了许多可怕的灾难。首先，他们驱逐我们，尽管我们孤身一人，被众人追赶，随时会被杀害，我们仍然在这样的时刻守了我们的节日。每个曾是我们中任何人连续受苦之地的地方，都成了我们集会的场所——田野、荒漠、船只、客栈、监狱，全都一样。然而，最欢乐的节日是由那些完美的殉道者庆祝的，他们已在天上参加宴席。此后，战争和饥荒突然降临我们。这些灾难我们与外邦人一同承受。但我们也独自承受了他们加给我们的那些祸害，同时我们也在他们彼此所做或所受的事上承担我们的份；而我们在基督所赐给我们的平安中深深喜乐，这平安是唯独赐给我们的。\par

3. 我们和他们一起享受了非常短暂的安息之后，这场瘟疫接着袭击了我们——这灾难对他们来说确实比一切可怕的东西更可怕，比任何其他种类的苦难更难以忍受；正如他们自己的某位作家所说，这是一种唯独战胜一切希望的灾难。但对我们来说却不是这样；它与其他灾难一样，成了我们训练和考验的工具。因为它并没有远离我们，尽管它在外邦人中传播得最为猛烈。\par

\Emph{他接着按顺序补充了这些话：}确实，我们许多弟兄，因着他们极大的爱心和兄弟情谊，他们不顾惜自己，而是彼此扶持，探望病人而不顾自己的危险，殷勤服侍他们，为了在基督内医治他们而不断喜乐地与他们一同死去，自己背负从别人而来的痛苦，将邻人的疾病引到自己身上，甘心将周围人的苦难担在自己身上。许多这样医好了别人疾病、使他们恢复健康的人，自己却死了，将死者身上的死亡转移到自己的身体上。那惯常的说法，在其他时候似乎只是一种礼貌的称呼，那时他们却在实际行为中活出来了，如同\Emph{万物的渣滓}一样离世。是的，我们最好的弟兄就这样离世了，包括一些长老和执事，以及在平信徒中最有声望的人：因此，这种死亡方式，因其中所显示的卓越虔诚和坚定信德，似乎一点也不亚于殉道本身。\par

5. 他们用翻过来的手和怀抱抬起圣人们的遗体，合上他们的眼睛，闭上他们的嘴。他们一同抬着，体面地铺陈，紧贴着他们，拥抱他们，用洗涤和衣着妥善预备。然后不久他们自己也同样接受了这些服务，因为幸存者总是跟随先逝者。但在外邦人中，一切恰恰相反。他们推开任何开始生病的人，甚至远离最亲密的朋友，将病人半死不活地扔在公共道路上，任凭他们不被埋葬，当他们死时就极其轻蔑地对待，坚定地避免与死亡有任何形式的交流和接触；然而，尽管他们采取了多种预防措施，却不容易完全逃脱。\par

% structure: p0065 source=h2 target=section reason=relative-heading-rank; base=h2

\section{书信十三。致\Place{埃及}主教\Person{希拉克斯}。}

1. 然而，如果我难以与定居在偏远地区的人通过书信沟通，又有什么奇怪呢？何况我似乎连与自己推理、与自己的灵魂商议的能力都不足。因为对我来说，与那些犹如我自己的内脏、最亲密的同伴和弟兄——与我同心，且同属一个教会的人——书信往来是非常必要的。然而，却没有途径让我传递这些书信。的确，一个人不仅是要走出行省边界，而且要从东到西跨越，也比从这个\Place{亚历山大}到\Place{亚历山大}旅行更容易。因为在这个城市中，最中心的道路比那片广阔无人的旷野更加荒凉难行，那旷野\Place{以色列}人只在两代人中穿越过；而我们平坦无波的海港已成了那片海的写照，当时那海分开，两边像墙一样立起，\Place{埃及}人就在那条通道中被淹死。因为常常由于在其中发生的屠杀，它们看起来像\Place{红海}。流经城市的河流，有时看起来比无水的沙漠更干涸，比那旷野更干旱，在那旷野中\Place{以色列}人如此口渴，以至于他们不断向\Person{梅瑟}抱怨，而水由那独行奇事者的能力从峭壁岩石中为他们流出。有时，河水又泛滥高涨，淹没周围的一切土地、道路和田野，仿佛要再次给我们带来\Person{诺厄}时代的大洪水。\par

2. 但如今，它总是向前流淌，被鲜血、屠杀和人的溺水挣扎所污染，正如古时因\Person{法郎}的缘故，被\Person{梅瑟}变为血水，变得腐臭。还有什么液体能清洁水本身呢？水本身能清洁一切，却无法清洁自己。那既广阔又难以跨越的海洋，即使倾注其上，又怎能净化这苦海？甚至那条从\Place{伊甸}流出的巨河，即使它将它分成的四道源头都排入\Place{基红}这一条河道，又怎能洗去这些污秽？这被四面八方升起的毒气污染的空气，何时才能变得纯净？因为从大地发出这样的蒸汽，从海洋发出这样的气流，从河流发出这样的微风，从港口发出这样的恶臭雾气，以至于我们以为露水就是所有底层元素中腐烂尸体的污浊流体。然而，尽管如此，人们仍然惊讶，不明白这些不断的瘟疫从何而来，这些可怕的疾病从何而来，这些多种致命的打击从何而来，这大量而多样的对生命的毁灭从何而来，以及为什么这座大城所容纳的总人口（从幼童到耄耋老人）不再像从前那样多，从前仅它称为健壮老人的就那么多。但那时从四十岁到七十岁的人比现在更多，以至于现在即使将从十四岁到八十岁的人都加进公共粮食配给的名册上，也无法凑足那个数目。而那些表面上看最年轻的人现在也变得，仿佛与从前最年长的人同龄。然而，尽管他们看到人类在地球上不断减少和消耗，他们在自己人数不断增长和加速的消耗与消灭中，却毫无恐惧。\par

% structure: p0068 source=h2 target=section reason=relative-heading-rank; base=h2

\section{第十四封书信。选自他的\WorkTitle{第四封节日书信}。}

爱总是时刻警惕，并设法向甚至不愿接受它的人行善。而且，许多时候，一个人出于羞愧而退缩，并因不愿给别人添麻烦而拒绝接受善意服务，宁愿自己承担痛苦也不愿引起任何人的烦恼和焦虑，但那个充满爱的人会强求他忍受帮助，并让他接受他人的援助，尽管这可能像在承受一种委屈，从而做成一件非常大的善事，不是对他自己，而是对别人，因为他允许别人成为终结他所陷入的邪恶的代理人。\par

\SourceRecord{Translated by S.D.F. Salmond.；Ante-Nicene Fathers；Vol. 6.；Edited by Alexander Roberts, James Donaldson, and A. Cleveland Coxe.；Buffalo, NY: Christian Literature Publishing Co.,；1886.；<http://www.newadvent.org/fathers/0632.htm>.}
